% Осень26!B7; нумерация внутри столбца без пустых строк.
% Источник: лекция 7, §10 (строка 335)
\definitionticket{05}{Формулировка теоремы о локальной обратимости}
\label{definition:05}

\subsection*{Теорема о локальной обратимости (без доказательства); формулировка в терминах систем уравнений}

Это третья теорема «из жизни диффеоморфизмов»: первая --- о сохранении области, вторая --- о гладкости обратного отображения, третья --- о локальной обратимости.

\begin{theorem}

\emph{о локальной обратимости и № 6.} Пусть \(O\subset\mathbb R^m\) открыто, \(T\colon O\to\mathbb R^m\) класса \(C^1\), и в точке \(x_0\in O\) производный оператор невырожден: \[\det T'(x_0)\ne0.\] Тогда существует окрестность \(U\ni x_0\) такая, что сужение \(T\big|_U\) --- \textbf{диффеоморфизм} (биективное гладкое отображение с гладким обратным).

\end{theorem}

\textbf{Замечания.}

\begin{itemize}
\tightlist
\item
  Так как \(\det T'\) непрерывен (ведь \(T\in C^1\): непрерывная функция от непрерывных матричных элементов) и отличен от нуля в \(x_0\), он отличен от нуля и в целой окрестности \(x_0\): невырожденность «автоматически» локальна.
\item
  \emph{Где новость.} Гладкость в прямую сторону дана по условию; открытость образа и гладкость обратного (после биективности) --- предыдущие теоремы. Новость ровно одна: \textbf{существует окрестность, на которой \(T\) биективно}. «Вот здесь неожиданно на нас свалилось, что оно обратимо.»
\item
  \emph{Почему «на пальцах» оно должно быть верно.} В первом порядке \(T\) близко к линейному отображению с невырожденным оператором; линейное невырожденное отображение точки не склеивает, и «почти линейное» локально тоже не склеивает --- усиленный вариант леммы о почти локальной инъективности. Для доказательства надо «помедитировать над определением дифференцируемости, превратив его в оценку», после чего всё получается само. \textbf{Требует сверки.} \textbf{Доказательство лектором пропущено}: «я позволю себе эту теорему не доказывать... она не сильно отличается от того, что я только что делал, просто немного другие технические шаги. Хорошая новость --- без доказательства.»
\end{itemize}
